\section{Chord}

Chord est un protocole DHT (\emph{Distributed Hash Table}) qui r\'esout un probl\`eme unique : trouver rapidement le n\oe ud responsable d'une cl\'e dans un r\'eseau pair-\`a-pair dynamique.

\subsection{Principe de fonctionnement}

Chaque n\oe ud et chaque cl\'e re\c coivent un identifiant sur \(m\) bits (souvent via SHA-1), plac\'e sur un anneau modulo \(2^m\). Une cl\'e \(k\) est stock\'ee sur son \emph{successor} : le premier n\oe ud rencontr\'e dans le sens horaire dont l'identifiant est sup\'erieur ou \`a \(k\).

Cette r\`egle donne trois propri\'et\'es utiles :
\begin{itemize}
\item \textbf{R\'epartition} : le \emph{consistent hashing} distribue les cl\'es de mani\`ere quasi uniforme.
\item \textbf{Scalabilit\'e} : l'ajout/retrait d'un n\oe ud ne d\'eplace qu'une fraction des cl\'es.
\item \textbf{D\'ecentralisation} : aucun serveur central n'est requis.
\end{itemize}

\subsection{Routage efficace : Finger Table}

Pour \`eviter un parcours lin\'eaire de l'anneau, chaque n\oe ud maintient une \textbf{finger table} de \(m\) entr\'ees. L'entr\'ee \(i\) pointe vers le premier n\oe ud joignable \`a partir de \(n + 2^{i-1}\).

\[
\text{finger}[i].\text{start} = (n + 2^{i-1}) \bmod 2^m
\]

Cette structure permet des sauts exponentiels et un co\^ut moyen de recherche en \(O(\log N)\) messages.

\subsection{Algorithmes essentiels}

\textbf{1) Recherche de la cl\'e : \texttt{find\_successor(id)}}
\begin{itemize}
\item Si \texttt{id} est entre le n\oe ud courant et son \emph{successor}, on retourne ce \emph{successor}.
\item Sinon, on avance vers le \emph{closest\_preceding\_node(id)} dans la finger table et on r\'ep\`ete.
\end{itemize}

\textbf{2) Choix du prochain saut : \texttt{closest\_preceding\_node(id)}}
\begin{itemize}
\item Parcours de la finger table du plus grand saut au plus petit.
\item Retour du premier doigt situ\'e strictement avant \texttt{id} sur l'anneau.
\end{itemize}

\textbf{3) Arriv\'ee d'un n\oe ud : \texttt{join(n')}}
\begin{itemize}
\item Le nouveau n\oe ud contacte un n\oe ud connu \(n'\).
\item Il initialise son \emph{successor} via \texttt{find\_successor} puis construit sa finger table.
\end{itemize}

\textbf{4) Stabilisation : \texttt{stabilize()} et \texttt{notify()}}
\begin{itemize}
\item Ex\'ecution p\'eriodique pour corriger les pointeurs \emph{successor}/\emph{predecessor}.
\item Garantit la convergence vers un anneau coh\'erent apr\`es churn mod\'er\'e.
\end{itemize}

\textbf{5) Tol\'erance aux pannes : \texttt{successor list}}
\begin{itemize}
\item Chaque n\oe ud conserve plusieurs successeurs de secours.
\item En cas de panne du successeur principal, le routage continue via le suivant vivant.
\end{itemize}

\subsection{Complexit\'e (r\'esum\'e)}
\begin{itemize}
\item \textbf{Lookup} : \(O(\log N)\) messages.
\item \textbf{Maintenance d'un join/leave} : surco\^ut logarithmique (d\'epend de la politique de mise \`a jour).
\item \textbf{M\'emoire par n\oe ud} : \(O(\log N)\) entr\'ees de routage.
\end{itemize}

En pratique, Chord fournit une base simple et robuste pour le \emph{lookup} distribu\'e : routage logarithmique, maintenance locale p\'eriodique et bonne r\'esilience aux variations du r\'eseau.